\chapter{Croup}

\section*{Definition and Overview}
Croup is a common childhood infection of the upper airways, caused by viruses that inflame the voice box (larynx) and windpipe (trachea). It typically affects children between six months and three years and produces a distinctive barking cough, a hoarse voice, and noisy breathing on breathing in (stridor). Croup usually begins with a cold and worsens at night. Most cases are mild and can be managed at home, but some children need medical treatment with steroids, and a small number need hospital care, particularly if breathing becomes difficult.

\section*{Epidemiology}
Croup is very common, affecting around 3--5\% of children, with a peak between six months and three years of age. It is slightly more common in boys and occurs most often in autumn and early winter. Most children have only one episode, but a small number have recurrent croup. Parainfluenza viruses cause most cases, with other viruses such as RSV, influenza, and adenovirus responsible for the rest. Although common, severe croup requiring intensive care is rare.

\section*{Aetiology and Pathophysiology}
Croup is caused by viral infection of the upper airway.

\begin{itemize}
    \item \textbf{Parainfluenza viruses:} the most common cause
    \item \textbf{Other viruses:} RSV, influenza, adenovirus, and rhinovirus
    \item \textbf{Inflammation:} the virus infects the lining of the larynx and trachea, causing swelling
    \item \textbf{Narrowing:} the airway is narrowest just below the vocal cords, so even mild swelling narrows it significantly
    \item \textbf{Stridor:} turbulent airflow through the narrowed airway causes the characteristic noisy breathing
    \item \textbf{The barking cough:} caused by swelling of the vocal cords
    \item \textbf{Severity:} worse in younger children, whose airways are smaller
\end{itemize}

\section*{Clinical Features}
Croup often begins like a cold, with symptoms peaking at night.

\begin{itemize}
    \item A seal-like barking cough
    \item A hoarse voice
    \item Noisy breathing on breathing in (inspiratory stridor)
    \item A runny nose, mild fever, and sore throat
    \item Symptoms worse at night and when upset or crying
    \item Mild chest wall indrawing with breathing in more severe cases
    \item Generally a well child between episodes of coughing
    \item Symptoms that typically last three to four days
\end{itemize}

\section*{Differential Diagnosis}
Other causes of noisy breathing and cough in children are considered, especially when symptoms are atypical.

\begin{itemize}
    \item Epiglottitis, a rare but dangerous bacterial infection of the flap above the voice box, with drooling and a toxic-looking child
    \item Bacterial tracheitis, a serious bacterial infection of the windpipe
    \item Inhaled foreign body, causing sudden coughing or wheezing
    \item Asthma or bronchiolitis, with wheeze rather than stridor
    \item Anaphylaxis, with rapid onset of swelling and breathing difficulty
    \item A peritonsillar or retropharyngeal abscess
    \item Congenital airway narrowing, in recurrent or atypical cases
\end{itemize}

\section*{When to Seek Medical Care}
Most croup is managed at home, but medical review is needed for moderate symptoms or worsening breathing.

\textbf{See a doctor promptly if:}
\begin{itemize}
    \item The stridor is present at rest or is getting louder
    \item The child is working hard to breathe, with indrawing of the chest
    \item The child is unsettled, restless, or very distressed
    \item Symptoms are severe or recurrent
    \item The child is under six months old
\end{itemize}

\textbf{Call 000 (or local emergency services) for:}
\begin{itemize}
    \item Difficulty breathing, or a child who is struggling or gasping
    \item A bluish or pale colour, or very low oxygen
    \item Drowsiness, floppiness, or difficulty waking
    \item Drooling and an inability to swallow (possible epiglottitis)
    \item A silent chest with little or no breathing sound
\end{itemize}

\section*{Diagnosis and Investigations}
Diagnosis is clinical in typical cases, and tests are reserved for severe or atypical presentations.

\begin{itemize}
    \item \textbf{History and examination:} the barking cough, stridor, fever, and assessment of breathing effort
    \item \textbf{Oxygen monitoring:} pulse oximetry in moderate or severe cases
    \item \textbf{No routine tests:} X-rays and blood tests are not needed for typical croup
    \item \textbf{Neck X-ray:} in atypical cases, to exclude epiglottitis or a foreign body
    \item \textbf{Hospital assessment:} for children with severe or worsening symptoms
\end{itemize}

\section*{Management}
Treatment depends on severity.

\begin{itemize}
    \item \textbf{Mild croup at home:} comfort, fluids, and paracetamol or ibuprofen for fever; sitting the child upright and keeping them calm
    \item \textbf{Steroids:} a single dose of oral dexamethasone or prednisolone, which reduces airway swelling and the need for hospital care
    \item \textbf{Nebulised adrenaline:} for moderate to severe croup in hospital, giving rapid temporary relief
    \item \textbf{Oxygen:} for children with low oxygen levels
    \item \textbf{Hospital admission:} for severe croup, or when breathing difficulty does not respond to treatment
    \item \textbf{Observation:} children with severe croup are monitored closely, as symptoms can recur
    \item \textbf{Intensive care:} rarely needed, for airway obstruction
\end{itemize}

\section*{Complications}
\begin{itemize}
    \item Worsening airway obstruction with severe breathing difficulty
    \item Hypoxia (low oxygen) affecting the brain and other organs
    \item Dehydration from poor drinking
    \item Secondary bacterial infection, such as pneumonia or tracheitis
    \item Recurrent croup in some children
    \item Rarely, the need for intensive care and intubation
\end{itemize}

\section*{Prognosis and Outcomes}
Croup is usually a mild, self-limiting illness that settles within three to four days, although the cough can linger. Treatment with steroids markedly reduces the severity and duration of symptoms and the need for hospitalisation. Severe complications are rare in otherwise healthy children. Children with recurrent croup should be reviewed, as some have underlying airway sensitivity or, rarely, an anatomical problem.

\section*{Prevention and Self-Care}
\begin{itemize}
    \item Keep the child calm and upright during episodes; crying worsens the stridor
    \item Offer small, frequent drinks to prevent dehydration
    \item Treat fever with age-appropriate paracetamol or ibuprofen
    \item Do not use steam or old-fashioned croup tents, which are not recommended
    \item Keep the child away from smoke
    \item Practise good hand hygiene to reduce spread of respiratory viruses
    \item Keep immunisations up to date, including influenza vaccination
    \item Know the warning signs that require urgent review
\end{itemize}

\section*{Sources and Further Reading}
The following reputable sources provide further information for healthcare students and patients seeking reliable, current advice about croup.

\begin{itemize}
    \item Healthdirect Australia. \textit{Croup.} https://www.healthdirect.gov.au/
    \item Royal Children's Hospital Melbourne. \textit{Kids Health Info: croup.} https://www.rch.org.au/kidsinfo/
    \item Raising Children Network. \textit{Croup.} https://raisingchildren.net.au/
    \item NHS. \textit{Croup.} https://www.nhs.uk/conditions/croup/
    \item Sydney Children's Hospitals Network. \textit{Croup fact sheet.} https://www.schn.health.nsw.gov.au/
    \item Mayo Clinic. \textit{Croup: overview and treatment.} https://www.mayoclinic.org/diseases-conditions/croup/
\end{itemize}
